\documentclass[aspectratio=169]{beamer}

% Sprache und Schriften (mit XeLaTeX kompilieren)
\usepackage{fontspec}
\usepackage[ngerman]{babel}
\usepackage{amsmath}

\usetheme{default}
\setbeamertemplate{navigation symbols}{}
\setbeamertemplate{footline}[frame number]
\definecolor{akzent}{HTML}{6B5CE7}
\setbeamercolor{structure}{fg=akzent}

\title{Titel der Präsentation}
\subtitle{Untertitel}
\author{Vorname Nachname}
\institute{Institution}
\date{\today}

\begin{document}

\begin{frame}
  \titlepage
\end{frame}

\begin{frame}{Inhalt}
  \tableofcontents
\end{frame}

\section{Motivation}

\begin{frame}{Motivation}
  \begin{itemize}
    \item<1-> Erster Punkt.
    \item<2-> Zweiter Punkt, der danach erscheint.
  \end{itemize}
\end{frame}

\section{Ergebnisse}

\begin{frame}{Ein Ergebnis}
  \begin{block}{Wellengleichung}
    \[
      \frac{\partial^2 u}{\partial t^2} = c^2 \nabla^2 u
    \]
  \end{block}
\end{frame}

\begin{frame}
  \centering\Large Vielen Dank!
\end{frame}

\end{document}
